\chapter{Záver}
\label{chapter:conclusion}

\section{Čo práca prináša}

Práca vychádza z pozorovania, že empirické zistenia o vplyve umelej inteligencie na učenie si protirečia, a z tvrdenia, že rozpor je zdanlivý. Zjednocujúca hypotéza formulovaná v časti~\ref{sec:syntezahypotezy} znie, že nástroj takmer vždy zvyšuje výkon počas hodiny a subjektívny pocit naučenia, kým jeho účinok na trvalé osvojenie závisí od charakteru interakcie žiaka s ním. Práca prináša deväť vecí.

\paragraph{Pozorovaná miera namiesto sebahodnotenia.} Index epistemickej kvality promptov prevádza doteraz sebahodnotenú veličinu~--- mieru kognitívneho offloadingu~--- na veličinu vypočítanú z toho, čo žiak skutočne napísal. Taxonómia je ukotvená v rámci ICAP~\cite{chi2014icap}, takže jej usporiadanie nie je ľubovoľné, ale odvodené z teórie, ktorá predpovedá učebné výsledky. Automatický klasifikátor postavený na priehľadných lexikálnych príznakoch dosahuje na kotviacom súbore zhodu s ľudskými kódmi vyššiu než simulovaná zhoda dvoch ľudí navzájom, pričom všetky jeho chyby sú posunom o jedinú susednú úroveň.

\paragraph{Manipulácia, ktorá je overiteľná, nie iba opísaná.} Všetky tri podmienky obsluhuje jediný program (časť~\ref{sec:cvicne-rozhranie}). Tvrdenie, že sa líšia výhradne pedagogickou vrstvou, tak prestáva byť vetou v metodike a stáva sa vlastnosťou, ktorú možno otestovať~--- a testuje sa: podmienka $K$ odmieta asistenta na strane servera a odmietnutie zapisuje ako odchýlku od protokolu; podmienky $G$ a $A$ sa v tom, čo odošlú modelu, líšia na dvoch riadkoch a v ničom inom. Výstupný guard navyše meria, ako často systémový prompt zlyháva, pretože \emph{miera} jeho zlyhania je číslo, ktoré práca musí vykázať, nie detail, ktorý má zamlčať.

\paragraph{Chyba merania mediátora ako samostatná otázka.} Literatúra o mediácii predpokladá, že mediátor je pozorovaný. Tu nie je~--- je odhadnutý klasifikátorom, a chyba toho odhadu môže závisieť od podmienky, ktorú dizajn randomizuje, pretože podmienka $A$ žiakom priamo predpisuje jazykovú formu, ktorú klasifikátor číta. Časť~\ref{sec:diferencialna} odvodzuje, čo taká chyba urobí s každým z troch odhadov, a výsledok je nesamozrejmý: redukovaná forma je voči nej úplne odolná, inštrumentálny odhad je ňou utlmený, a mediačný odhad je najprv nafúknutý a potom utlmený tak, že sa obe skreslenia môžu presne zrušiť. Z toho plynie poradie, v akom sa mechanizmové výsledky uvádzajú, a rozsah zaslepeného auditu, ktorým sa zvyšková chyba odhadne.

\paragraph{Dizajn, ktorý izoluje pedagogiku od modelu.} Doterajšie porovnania adaptívnych a generických nástrojov zamieňajú model, rozhranie, obsah aj pedagogiku~--- rozdiel preto nemožno pripísať žiadnemu z nich. Navrhnutý dizajn drží model, úlohy, rozhranie, učiteľa aj čas konštantné a mení výhradne pedagogickú vrstvu. Ak sa podmienky $G$ a $A$ budú líšiť, je to dôkaz o účinku ohraničenia, nie o kvalite modelu.

\paragraph{Meranie ilúzie porozumenia.} Položková kalibrácia sebadôvery umožňuje oddeliť tri veci, ktoré sa bežne zamieňajú: nadhodnotenie, absolútnu presnosť odhadu a schopnosť rozlíšiť, ktorým vlastným odpovediam veriť. Simulácia naznačuje, že tieto tri sa pri práci s nástrojom nemusia meniť spolu, čo má priamy praktický dôsledok pre to, akú nápravu má zmysel navrhovať.

\paragraph{Ukazovateľ, ktorý sa nesmie pohnúť.} Dizajn obsahuje dve negatívne kontroly s vopred stanoveným rozhodovacím pravidlom (časť~\ref{sec:negativne-kontroly}). Bez nich je rozdiel medzi „žiaci sa naučili viac“ a „žiaci sa viac snažili na teste, ktorý nasledoval po hodine, ktorá ich bavila“ z primárneho kontrastu nezistiteľný. S nimi je: motivácia, zvyšná kapacita aj nálada predpovedajú posun na látke, ktorú hodina nepokryla, kým učenie ho nepredpovedá. Práca zároveň vykazuje, akú hranicu kontrola dosiahne a akú nie, pretože negatívna kontrola bez uvedenej hranice funguje ako ubezpečenie, nie ako dôkaz.

\paragraph{Rozlíšenie toho, čo dizajn preukazuje, od toho, čo predpokladá.} Randomizácia prideľuje podmienku, nie kvalitu promptov. Práca preto zaraďuje do konfirmačnej rodiny len tvrdenia, ktoré sú randomizáciou identifikované, a mechanizmus testuje trianguláciou troch odhadov pod tromi rôznymi predpokladmi, doplnenou o analýzu citlivosti na nepozorované zmätenie. Toto rozlíšenie sa v literatúre o umelej inteligencii vo vzdelávaní robí zriedka, hoci práve mechanizmové tvrdenia sú tie, ktoré sa najčastejšie prenášajú do odporúčaní pre prax.

\paragraph{Analytický plán overený pred zberom.} Celý štatistický plán bol implementovaný a spustený na dátach zo známeho generujúceho modelu. Postup odhalil deväť chýb návrhu, ktoré by sa inak prejavili až po zbere, keď je náprava nemožná~--- medzi nimi singularitu v kódovaní prenosu medzi periódami, chýbajúce prepojenie testových foriem kotviacimi položkami a prítomnosť neidentifikovateľnej hypotézy v konfirmačnej rodine.

\paragraph{Reprodukovateľnosť ako vlastnosť, nie ako disciplína.} Žiadne číslo v kapitolách~\ref{chapter:metodika} a~\ref{chapter:vysledky} nie je napísané rukou; všetkých \repoMacros{} vzniká generovaním a text sa na ne odvoláva menom. Rovnaké pravidlo platí pre čísla, ktoré práca uvádza o sebe samej. Doplnkovo beží \repoChecks{} konzistenčných kontrol a \repoTests{} automatických testov v \repoSuites{} sadách. Rozdelenie medzi nimi je podstatné a v literatúre sa zamieňa: kontroly overujú, či sa text zhoduje s výstupmi analýzy, testy overujú, či je analýza správna. Prvé sú bezcenné bez druhých~--- zle naprogramovaný odhad by prešiel všetkými kontrolami, lebo práca by jeho zlé číslo vytlačila verne.

\section{Inžiniersky a spoločenský dosah}

Technický artefakt prešiel jednou úplnou iteráciou, ktorej výsledok bol negatívny a práve preto hodnotný: overovacia kostra postavená na knižnici LittleProxy nedokáže zachytiť komunikáciu moderných konverzačných rozhraní. Zistenie prišlo v čase, keď bola zmena lacná. Revidovaná architektúra izoluje záchytnú vrstvu za súborovým rozhraním, takže jej budúca výmena~--- ktorá nepochybne príde~--- nezasiahne logiku, ktorá je predmetom práce.

Druhé inžinierske zistenie je všeobecnejšie a prenosné mimo tejto práce. Rovnaká špecifikácia implementovaná dvakrát sa rozíde, a rozíde sa ticho: opravená chyba v jednej implementácii zostala v druhej, pričom výsledkom by neboli chybové hlásenia, ale mierne iné čísla (časť~\ref{sec:dve-implementacie}). Obrana, ktorá funguje, nie je zjednotenie kódu~--- to často nie je možné~--- ale zafixovanie špecifikácie ako \emph{dát}, na ktorých sa obe implementácie musia zhodnúť. Podobne poučná je trieda chýb, ktorú prekladač ani kontrola krížových odkazov neodhalí: príkaz \LaTeX{}u, ktorý prišiel o spätné lomítko, sa vytlačí ako text a nechýba, takže kontrola, že každý odkaz vedie na cieľ, ho nevidí. Kontrolovať treba príčinu, nie následok.

Spoločenský dosah je viazaný na hypotézu H7. Ak platí, že poškodenie generickým nástrojom rastie s klesajúcimi vstupnými vedomosťami, potom plošné sprístupnenie nástroja bez pedagogického ohraničenia zväčšuje rozdiely medzi žiakmi~--- hoci argumentom pre jeho zavedenie býva pravý opak. Praktickým výstupom práce nie je odporúčanie „používať“ alebo „nepoužívať“, ale rozlíšenie medzi tým, čo sa oplatí zaviesť, a tým, čo je len pohodlné.

Druhý dosah sa týka toho, ako sa o zavádzaní rozhoduje. Ak sa potvrdí disociácia medzi pocitom naučenia a trvalým osvojením, potom dotazník spokojnosti žiakov~--- najdostupnejší nástroj, ktorý škola má~--- nie je pri umelej inteligencii len nedostatočný, ale systematicky zavádzajúci: uprednostní horšiu možnosť.

\section{Obmedzenia}

Obmedzenia uvádzame otvorene, vrátane tých, ktoré ošetriť nevieme.

\begin{description}[leftmargin=0.6cm, nosep]
    \item[\textbf{Jediná doména.}] Tri témy z algoritmizácie nedovoľujú tvrdiť nič o dejepise či o cudzom jazyku. Voľba domény bola vedená internou validitou~--- potrebou objektívne skórovateľného testu a izomorfných tém~--- a jej cenou je obmedzená zovšeobecniteľnosť.
    \item[\textbf{Zámerný výber škôl.}] Školy nie sú vybrané náhodne, pretože to v školskom výskume nie je možné. Stratifikácia podľa typu školy a lokality riziko zmierňuje, neodstraňuje ho. Náhodný výber je nemožný, ale \emph{zmerať}, kde vzorka leží, možné je: kotviace úlohy iBoboru so zverejnenou národnou úspešnosťou umiestnia vzorku na národnú škálu so smerodajnou chybou \nanSe{} logitu (časť~\ref{sec:narodne-ukotvenie}). Reprezentatívnou ju to nerobí; mení to nepreukázateľný predpoklad na vykázané číslo, pri ktorom čitateľ vidí smer a veľkosť odchýlky.
    \item[\textbf{Neviditeľné použitie.}] Nástroj zachytí len prevádzku zo školského zariadenia. Meranie miery pokrytia umožňuje rozsah problému kvantifikovať, nie ho odstrániť.
    \item[\textbf{Krátky odstup.}] Tri týždne sú viac než vo väčšine porovnateľných štúdií, ale stále málo na tvrdenie o trvalom osvojení v zmysle školského roka.
    \item[\textbf{Útlm chybou merania.}] Testy s dvanástimi až šestnástimi položkami podhodnocujú účinky približne o pätinu. Dlhší test by nevošiel do vyučovacej hodiny.
    \item[\textbf{Jedna verzia modelu.}] Zistenia o konkrétnom modeli zastarávajú rýchlo. Hlavné tvrdenie práce je preto formulované o pedagogickom ohraničení, nie o modeli.
    \item[\textbf{Neoverená pedagogická vrstva.}] Adaptívna vrstva je vlastným návrhom, nie validovaným produktom. Jej správanie sa overuje sadou štyridsiatich útočných dopytov s požiadavkou na nulový únik, strojovým guardom nad každou odpoveďou a kontrolou vzorky odpovedí. To je viac než opis v metodike, ale stále to nie je nezávislá validácia: všetky tri kontroly navrhol ten istý človek ako vrstvu samotnú.
    \item[\textbf{Klasifikátor je kalibrovaný na kotviacom súbore, nie na triede.}] Útok formulkou je uzavretý na 90~kotviacich promptoch (časť~\ref{sec:dominancia}) a zhoda dvoch nezávislých implementácií je testovaná na 129~prípadoch. Ani jedno nie je dôkaz, že v skutočnej triede nezostane diferenciálna chyba; presne preto je v pláne zaslepený audit, ktorý ju odhadne priamo na dátach štúdie (časť~\ref{sec:audit}). Zostáva jeden prípad, ktorý ani audit nezachytí: keby sa ľudskí kódovatelia mýlili rovnakým smerom ako stroj.
    \item[\textbf{Mechanizmus nie je identifikovaný dizajnom.}] Randomizácia prideľuje podmienku, nie kvalitu promptov. Tvrdenie, že účinok prechádza cez PEQI, zostáva podmienené predpokladmi, ktorých porušenie vieme kvantifikovať, ale nie vylúčiť. Na simulovaných dátach vychádza prah citlivosti $\rho^\ast$ nízky, takže mechanizmové tvrdenie by sa v takom prípade uvádzalo len ako otvorená otázka.
    \item[\textbf{Citlivosť na výpadok.}] Hlavný záver prežije hodnoverný vzorec výpadku, nie však maximálne nepriaznivý. Udržanie účasti na odloženom teste je preto pre túto štúdiu dôležitejšie než rozšírenie vzorky~--- čo je zistenie, ktoré mení priority pri organizácii zberu.
\end{description}

\section{Ďalší postup}

Bezprostredným krokom je dokončenie výmeny záchytnej vrstvy, pilotné overenie na jednej triede a získanie schválení, ktoré tvoria kritickú cestu celého projektu (príloha~\ref{chapter:Work-Schedule}). Zber prebehne podľa plánu opísaného v časti~\ref{sec:priebeh} a analýza podľa plánu, ktorý je už uzavretý a nebude sa meniť.

Z hľadiska ďalšieho výskumu sa ponúkajú štyri smery, ktoré z práce priamo vyplývajú. \emph{Nezávislá validácia indexu}~--- prototyp herného nástroja s uzavretými možnosťami opísaný v časti~\ref{sec:temy} poskytuje skóre kritického posúdenia výstupov nástroja, ktoré je odvodené inak než PEQI; ak obe veličiny merajú to isté, musia spolu korelovať, a ak nekorelujú, je to informácia o indexe, nie o hre. \emph{Prenos indexu do inej domény}~--- ak je taxonómia promptov naozaj ukotvená v rámci ICAP a nie v algoritmike, mala by fungovať aj na esejách či na cudzom jazyku; to je testovateľné. \emph{Intervencia na úrovni promptu}~--- ak epistemická kvalita promptov predpovedá trvalé osvojenie, potom priame učenie žiakov, ako sa nástroja pýtať, by malo mať merateľný účinok, a to lacnejšie než budovanie adaptívnej vrstvy. \emph{Dlhodobé sledovanie}~--- Kosmyna a kolektív~\cite{kosmyna2025brain} zistili, že rozdiely sa v priebehu štyroch mesiacov neznižovali, ale prehlbovali; či to platí aj pre školské učivo, je otvorená otázka s vážnymi dôsledkami.

Nakoniec je potrebné povedať, čo práca urobiť nemôže. Nerozhodne, či majú školy nástroje umelej inteligencie zavádzať~--- to je rozhodnutie, ktoré závisí aj od hodnôt, nákladov a cieľov mimo dosahu jedného experimentu. Dokáže však poskytnúť to, čo v tomto rozhodovaní na Slovensku zatiaľ chýba úplne: domáce dáta a merací nástroj, ktorým sa dá overiť, či zavedená forma funguje, alebo sa len páči.
